\chapter{Fundamentals of Visual-Manual Description and Identification of Soils}

\section{Historical Development of Visual-Manual Soil Identification}

\subsection{Origins of Empirical Soil Knowledge in Ancient Civilizations}

\subsubsection{Early Soil Classification Practices in China (circa 2000 BCE)}

\subsubsection{Greco-Roman Agricultural and Constructive Soil Observations}

\subsubsection{Medieval and Renaissance Contributions to Soil Description}

\subsection{Emergence of Scientific Soil Identification in the Modern Era}

\subsubsection{Nineteenth-Century Geological and Agricultural Roots}

\subsubsection{Early Geotechnical Engineering Practices and Field Identification}

\subsubsection{Karl Terzaghi and the Foundations of Systematic Soil Mechanics}

\subsection{Development of Standardized Engineering Classification Systems}

\subsubsection{Casagrande's Airfield Classification System (1942--1948)}

\subsubsection{Establishment of the Unified Soil Classification System (USCS, 1952)}

\subsubsection{Integration of Visual-Manual Procedures into ASTM Standards}

\subsubsection{Evolution of ASTM D2488: From 1966 to the 2017 Edition}

\section{Scope, Significance, and Epistemological Basis of Visual-Manual Procedures}

\subsection{Conceptual Distinction Between Description and Identification}

\subsection{Role of Visual-Manual Procedures in the Geotechnical Investigation Workflow}

\subsection{Relationship Between ASTM D2488 and ASTM D2487}

\subsection{Applicability to Naturally Occurring and Anthropogenic Materials}

\subsection{Limitations and Professional Judgment in Visual-Manual Practice}

\section{Comparative Analysis with Descriptive Methodologies in Allied Disciplines}

\subsection{Geological Field Description of Soils and Sediments}

\subsubsection{Sedimentological Description Protocols}

\subsubsection{Geomorphological and Quaternary Geology Approaches}

\subsection{Pedological Description Methods}

\subsubsection{FAO Guidelines for Soil Description}

\subsubsection{USDA Soil Survey Manual and Horizon-Based Description}

\subsubsection{ISO 11259 and ISO 25177 for Pedological Purposes}

\subsection{Agricultural and Agropedological Visual Soil Assessment}

\subsubsection{Visual Soil Assessment (VSA) Methods}

\subsubsection{Soil Texture by Feel: Agronomic Versus Engineering Criteria}

\subsection{Environmental Science and Contaminated Land Assessment}

\subsubsection{Field Identification of Contaminated Soils}

\subsubsection{Organoleptic Indicators in Environmental Site Assessment}

\subsection{Comparative Synthesis: Commonalities and Divergences Across Disciplines}

\section{International Standards for Visual-Manual Soil Description: A Comparative Review}

\subsection{ISO 14688-1:2017 (Geotechnical Investigation and Testing -- Identification and Description of Soil)}

\subsubsection{Scope and General Principles}

\subsubsection{Material Characteristics and Mass Characteristics}

\subsubsection{Field Hand Tests and Descriptive Vocabulary}

\subsection{ISO 14688-2:2017 (Classification of Soil for Engineering Purposes)}

\subsection{Eurocode 7 (EN 1997-2) and Its Relationship to Soil Description Standards}

\subsection{British Standard BS 5930:2015 (Code of Practice for Ground Investigations)}

\subsubsection{Historical Evolution: CP2001, BS 5930:1981, BS 5930:1999, BS 5930:2015}

\subsubsection{Soil Description Procedures and Terminology}

\subsubsection{Comparison of Consistency and Strength Descriptors with ASTM D2488}

\subsection{Chinese Standard GB/T 50145-2007 (Standard for Engineering Classification of Soil)}

\subsubsection{Structural Organization and Classification Criteria}

\subsubsection{Field Identification Procedures and Descriptive Terms}

\subsubsection{Comparison with USCS-Based Approaches}

\subsection{Australian Standard AS 1726 (Geotechnical Site Investigations)}

\subsubsection{Scope and Structural Framework}

\subsubsection{Description Criteria and Terminology}

\subsection{Japanese Geotechnical Standards: JGS and JIS Frameworks}

\subsection{Ibero-American and Latin American Normative Frameworks}

\subsubsection{IRAM Standards (Argentina)}

\subsubsection{ABNT NBR Standards (Brazil)}

\subsubsection{Peruvian Standards: NTP and References to ASTM D2488}

\subsection{Comparative Analysis of Key Descriptive Parameters Across International Standards}

\subsubsection{Grain Size Classification Boundaries}

\subsubsection{Plasticity and Fine-Grained Soil Identification Criteria}

\subsubsection{Color, Structure, and Macrofabric Descriptors}

\subsubsection{Consistency and Strength Terminology}

\section{Fundamental Principles Underlying Visual-Manual Identification Criteria}

\subsection{Principles of Grain Size Estimation Without Laboratory Equipment}

\subsubsection{Perception Thresholds and Visual Grain Size Limits}

\subsubsection{Weight-Based Versus Volume-Based Percentage Estimates}

\subsubsection{Minimum Sample Mass Requirements and Their Scientific Basis}

\subsection{Particle Angularity and Shape: Morphological Basis and Engineering Relevance}

\subsubsection{Definitions and Descriptive Criteria for Angularity}

\subsubsection{Definitions and Criteria for Particle Shape (Flat, Elongated)}

\subsubsection{Influence of Angularity and Shape on Shear Strength and Compressibility}

\subsection{Soil Color: The Munsell System and Its Scientific Rationale}

\subsubsection{History and Development of the Munsell Color System}

\subsubsection{Hue, Value, and Chroma: Physical and Perceptual Foundations}

\subsubsection{Munsell Soil Color Charts: Structure and Practical Use}

\subsubsection{Pedogenic and Geotechnical Information Encoded in Soil Color}

\subsubsection{Effect of Moisture Content on Color and Standardization Protocol}

\subsection{Odor as a Diagnostic Property}

\subsubsection{Organic Decomposition and Volatile Compound Indicators}

\subsubsection{Contaminant Identification Through Olfactory Assessment}

\subsection{Moisture Condition: Assessment and Engineering Significance}

\subsubsection{Criteria for Dry, Moist, and Wet Classification}

\subsubsection{Relationship Between Moisture State and Field Behavior}

\subsection{Reaction with Hydrochloric Acid: Chemical Basis and Diagnostic Use}

\subsubsection{Carbonate Chemistry and the HCl Effervescence Test}

\subsubsection{Safety Protocols for Field Use of Dilute HCl}

\subsubsection{Engineering Implications of Carbonate-Rich Soils}

\subsection{Dry Strength Test: Basis in Soil Mineralogy and Plasticity}

\subsubsection{Preparation Protocol and Specimen Standardization}

\subsubsection{Correlation Between Dry Strength, Clay Mineralogy, and Atterberg Limits}

\subsection{Dilatancy Test: Physical Mechanism and Diagnostic Significance}

\subsubsection{Pore Pressure Response to Shaking and Squeezing}

\subsubsection{Dilatancy as an Indicator of Soil Type and Gradation}

\subsection{Toughness Test: Thread Rolling and Plasticity Assessment}

\subsubsection{Relationship Between Thread Toughness and the Plasticity Index}

\subsubsection{Thread Rolling Procedure and Plastic Limit Approximation}

\subsection{Plasticity: Casagrande's A-Line and Its Visual Counterpart}

\subsubsection{Atterberg Limits: Definition and Engineering Context}

\subsubsection{Estimating Plasticity Through Combined Manual Tests}

\subsection{Consistency of Intact Fine-Grained Soils}

\subsubsection{Undrained Shear Strength and Its Relationship to Consistency Classes}

\subsubsection{Field Assessment Using Thumb Penetration Test}

\subsection{Cementation: Mechanisms and Field Identification}

\subsubsection{Common Cementing Agents in Soils (Carbonates, Iron Oxides, Silica)}

\subsubsection{Finger Pressure Test and Strength Descriptors}

\subsection{Structure and Macrofabric of Intact Soils}

\subsubsection{Genetic Origin of Soil Structure: Depositional and Post-Depositional Processes}

\subsubsection{Description Categories: Stratified, Laminated, Fissured, Slickensided, Lensed, Blocky, Homogeneous}

\subsection{Hardness of Coarse Particles}

\subsubsection{Hammer Test Procedure and Interpretation}

\subsubsection{Mineralogical Basis for Particle Hardness Variability}

\section{Standard Order of Description and the Single-Paragraph Approach}

\subsection{Rationale for Standardized Descriptive Terminology and Sequence}

\subsection{ASTM D2488 Standard Word Order and Checklist}

\subsection{Group Symbol and Group Name Assignment}

\subsubsection{Identification of Coarse-Grained Soils: Gravels and Sands}

\subsubsection{Identification of Fine-Grained Soils: Silts and Clays}

\subsubsection{Identification of Highly Organic Soils: Peat (PT)}

\subsubsection{Borderline Symbols and Dual Classification Notation}

\subsection{In-Place Conditions: Consistency, Structure, Moisture, and Color of Intact Samples}

\subsection{Geologic Interpretation and Local Nomenclature}

\subsection{Additional Comments and Supplementary Observations}

\section{Practical Identification Procedures for Coarse-Grained Soils}

\subsection{Preliminary Assessment: Gravel Versus Sand Dominance}

\subsection{Identification of Clean Gravels and Sands}

\subsubsection{Well-Graded Soils (GW, SW)}

\subsubsection{Poorly Graded Soils (GP, SP)}

\subsection{Identification of Soils with Fines}

\subsubsection{Silty Gravels and Sands (GM, SM)}

\subsubsection{Clayey Gravels and Sands (GC, SC)}

\subsection{Dual Symbol Assignment for Intermediate Fines Content}

\subsection{Incorporation of Cobbles and Boulders in Description}

\section{Practical Identification Procedures for Fine-Grained Soils}

\subsection{Specimen Preparation and Representative Sample Selection}

\subsection{Execution of Dry Strength, Dilatancy, and Toughness Tests}

\subsection{Identification of Inorganic Fine-Grained Soils (CL, CH, ML, MH)}

\subsubsection{Lean Clay (CL) and Fat Clay (CH)}

\subsubsection{Silt (ML) and Elastic Silt (MH)}

\subsubsection{Distinguishing MH from CL: Field Indicators}

\subsection{Identification of Organic Fine-Grained Soils (OL, OH)}

\subsubsection{Color Change upon Air Drying}

\subsubsection{Organic Odor and Spongy Thread Behavior}

\subsection{Identification of Peat (PT)}

\subsection{Modifiers for Sand and Gravel Content in Fine-Grained Soils}

\section{Extended and Complementary Field Identification Procedures}

\subsection{Dispersion-Sedimentation Test}

\subsubsection{Procedure: Jar or Test Tube Method}

\subsubsection{Interpretation of Settling Times by Particle Size}

\subsubsection{Limitations for High-Plasticity Clays and Flocculated Suspensions}

\subsection{Tooth Test (Bite Test)}

\subsubsection{Physical Basis for Gritty Versus Smooth Sensation}

\subsubsection{Lower Bound of Visible Grain Size and Its Relationship to 0.075 mm}

\subsection{Sphericity and Roundness Estimation}

\subsubsection{Definitions and Conceptual Distinction from Angularity}

\subsubsection{Visual Reference Charts for Sphericity and Roundness Estimation}

\subsubsection{Influence on Minimum and Maximum Void Ratios}

\subsection{Behavior in Air and Water Test}

\subsubsection{Immersion Test: Disintegration of Silt Versus Cohesion of Clay}

\subsubsection{Drying Behavior on Smooth Surfaces: Clay Adhesion Versus Silt Dusting}

\subsubsection{Engineering Interpretation of Clay Versus Silt Behavioral Indicators}

\subsection{Particle Surface Coatings and Mineralogical Indicators}

\subsubsection{Presence of Mica, Gypsum, and Other Accessory Minerals}

\subsubsection{Iron Oxide and Organic Coatings on Coarse-Grained Particles}

\subsection{Root Presence, Root Holes, and Biogenic Features}

\subsection{Caving and Drilling Resistance as Indirect Soil Indicators}

\section{Documentation, Reporting, and Field Recording Standards}

\subsection{Field Log Formats and Standardized Description Forms}

\subsection{Photographic Documentation of the Site and Surrounding Environment}

\subsection{Photographic Documentation of Soil Samples}

\subsubsection{Sample Area Preparation and Reference Scale}

\subsubsection{Munsell Color Scale Integration in Field Photography}

\subsection{Sampling Context: Origin, Depth, and Method of Extraction}

\subsubsection{Borehole, Test Pit, and Auger Sample Identification}

\subsubsection{Disturbed Versus Undisturbed Sample Considerations}

\subsection{Abbreviated Classification Symbols for Graphical Logs and Databases}

\subsection{Geologic Interpretation and Supplementary Contextual Information}

\section{Application of the Visual-Manual System to Non-Soil Materials}

\subsection{Weakly Cemented Rocks and Transitional Materials}

\subsubsection{Shale, Claystone, Siltstone, and Mudstone}

\subsubsection{Slaking Behavior and Post-Slaking Identification}

\subsection{Anthropogenic and Industrial Materials}

\subsubsection{Shells, Crushed Rock, and Slag}

\subsubsection{Recycled Construction Materials and Fill}

\subsection{Frozen Soils: Interface with ASTM D4083}

\section{Sources of Error, Variability, and Quality Assurance in Visual-Manual Practice}

\subsection{Observer Variability and Subjectivity in Sensory Assessment}

\subsection{Sample Disturbance and Its Effect on Descriptive Accuracy}

\subsection{Experience, Training, and Calibration of Field Personnel}

\subsubsection{Learning Through Comparison with Laboratory Results}

\subsubsection{Role of Supervised Practice with Experienced Professionals}

\subsection{Repeatability and Reproducibility of Manual Tests}

\subsection{Integration with Laboratory Classification: Verification Protocols}

\section{Visual-Manual Description in the Context of Geotechnical Site Investigation}

\subsection{Borehole Logging and Stratigraphic Profiling}

\subsection{Test Pit and Trench Logging}

\subsection{Sample Grouping and Optimization of Laboratory Testing Programs}

\subsection{Correlation Between Visual Description and In Situ Testing Results}

\subsubsection{Standard Penetration Test (SPT) Correlations}

\subsubsection{Cone Penetration Test (CPT) Soil Behavior Type Classification}

\subsection{Visual-Manual Procedures in Geotechnical Risk Assessment}

\section{Emerging Trends and Future Perspectives in Field Soil Identification}

\subsection{Digital Color Measurement as a Complement to Munsell Charts}

\subsection{Machine Learning and Artificial Intelligence in Automated Soil Description}

\subsection{Hyperspectral and Remote Sensing Technologies for Field Identification}

\subsection{Standardization Challenges in a Globalized Geotechnical Practice}

\subsection{Toward a Universal Visual-Manual Framework: Prospects for International Harmonization}
